\chapter{Glossaire métier}
\label{annexe:glossaire}

Glossaire complet dans \filepath{docs/architecture/topographie_donnees.md}
\S{}1 (C2) ; repris ici in extenso pour que le rapport reste
autosuffisant. Chaque terme est une cible de renvoi (voir par exemple
\hyperref[gloss:scd2]{SCD2} au chapitre~\ref{chap:bloc3}) --
cliquable dans le PDF.

\begin{longtable}{@{}>{\bfseries\raggedright}p{3.4cm} p{12.2cm}@{}}
\toprule
Terme & Définition \\
\midrule
\endfirsthead
\toprule
Terme & Définition \\
\midrule
\endhead
\bottomrule
\endfoot

\label{gloss:3pl}3PL (\textit{Third-Party Logistics}) &
Prestataire logistique externalisé gérant entreposage, préparation de
commandes et transport pour le compte de tiers (les clients d'Omega
Logistics). \\

\label{gloss:wms}WMS (\textit{Warehouse Management System}) &
Système de gestion d'entrepôt. Chez Omega Logistics : \textbf{FluxPro},
qui pilote commandes, stocks, produits et expéditions. \\

\label{gloss:tms}TMS (\textit{Transport Management System}) &
Système de gestion du transport. Chez Omega Logistics :
\textbf{TransFlow}, qui organise tournées et livraisons de la flotte. \\

\label{gloss:entrepot}Entrepôt &
Site d'exploitation d'Omega Logistics. Trois sites : Lyon
(\texttt{OMG-LYO}), Lille (\texttt{OMG-LIL}), Marseille
(\texttt{OMG-MAR}). \\

\label{gloss:client}Client &
Donneur d'ordre d'Omega Logistics. Trois clients fictifs : NordDrive
(pièces auto), FreshMarket (grande distribution alimentaire), MedioTex
(textile). \\

\label{gloss:commande}Commande &
Demande d'expédition de produits passée par un client pour un entrepôt
donné, composée d'une ou plusieurs lignes de commande. \\

\label{gloss:expedition}Expédition &
Envoi physique résultant d'une commande, confié à un transporteur,
identifié par un numéro de suivi (\texttt{tracking\_number}). \\

\label{gloss:stock}Stock &
Quantité disponible d'un produit dans un entrepôt à une date de mise à
jour donnée. \\

\label{gloss:sku}SKU (\textit{Stock Keeping Unit}) &
Référence unique d'un produit dans le référentiel FluxPro. \\

\label{gloss:tournee}Tournée &
Ensemble de livraisons affectées à un chauffeur et un véhicule, pour
une date donnée, organisé par TransFlow. \\

\label{gloss:livraison}Livraison &
Remise d'un colis (une expédition) à une adresse, rattachée à une
tournée, avec un statut et des horaires estimé/réel. \\

\label{gloss:transporteur}Transporteur &
Société assurant le transport physique des expéditions (ex.
RapidFret). Distinct d'Omega Logistics, dont TransFlow orchestre les
tournées. \\

\label{gloss:sla}SLA (\textit{Service Level Agreement}) &
Accord de niveau de service définissant la qualité de prestation
attendue. \\

\label{gloss:rfid}RFID (\textit{Radio Frequency Identification}) &
Technologie d'étiquetage/scan sans contact utilisée pour tracer les
palettes en entrepôt. \\

\label{gloss:iot}IoT (\textit{Internet of Things}) &
Capteurs connectés en entrepôt (température, comptage caméra) et sur
la flotte (géolocalisation) -- bloc 4 (data lake). \\

\label{gloss:sse}SSE (\textit{Server-Sent Events}) &
Protocole HTTP de flux continu à sens unique serveur~$\rightarrow$~client
(une connexion ouverte, le serveur pousse des évènements). Utilisé par
\texttt{/api/stream/capteurs} pour diffuser les mesures capteurs en
temps réel (bloc 4). \\

\label{gloss:staging}Staging &
Base de données de travail intermédiaire, cible de la collecte du
bloc 2, alimentant ensuite l'entrepôt de données (bloc 3). À ne pas
confondre avec la zone \texttt{staging} du data lake (bloc 4), qui
porte le même nom par convention Kimball/data lake mais désigne une
brique de stockage différente. \\

\label{gloss:scd1}SCD1 (\textit{Slowly Changing Dimension}, type 1) &
Historisation Kimball où un changement \textbf{écrase} la valeur
précédente, sans garder de trace. Cas implicite dans OMEGA BI :
\texttt{Dim\_Client.code} n'est jamais versionné. \\

\label{gloss:scd2}SCD2 (\textit{Slowly Changing Dimension}, type 2) &
Historisation Kimball où un changement produit une \textbf{nouvelle
ligne} (nouvelle clé de substitution, \texttt{valid\_from}/
\texttt{valid\_to}/\texttt{is\_current}), en conservant l'ancienne.
Implémentée sur \texttt{Dim\_Client.nom}/\texttt{secteur} en C17 --
voir chapitre~\ref{chap:bloc3}. \\

\label{gloss:s3}S3 (\textit{Simple Storage Service}) &
Service de stockage objet d'AWS ; son API HTTP est devenue un standard
de fait, repris par des solutions non-AWS comme \textbf{MinIO}
(bloc 4) pour rester compatibles avec les mêmes outils sans dépendre
d'AWS. \\

\label{gloss:bucket}Bucket &
Conteneur de premier niveau d'un stockage objet S3/MinIO
(grossièrement, un répertoire racine). OMEGA LAKE utilise un
\textbf{bucket unique}, \texttt{omega-lake} -- les zones
\texttt{raw}/\texttt{staging}/\texttt{curated} sont des préfixes à
l'intérieur de ce bucket, pas des buckets séparés. \\

\label{gloss:zone-lake}Zone (data lake) &
Une étape de transformation à l'intérieur du bucket unique
(\texttt{raw}/\texttt{staging}/\texttt{curated}) -- un préfixe de
chemin, pas un bucket séparé. \\

\label{gloss:flux-lake}Flux (data lake) &
Une source de données IoT distincte, avec son propre schéma (ex.
\texttt{capteurs\_temperature}, \texttt{rfid\_scans}) -- un
sous-préfixe à l'intérieur de chaque zone. 5 flux au total. \\

\label{gloss:httpfs}httpfs &
Extension DuckDB qui sait lire/écrire des fichiers via HTTP(S) ou
\texttt{s3://}, y compris contre un stockage S3-compatible non-AWS
comme MinIO -- nécessite alors une configuration explicite d'endpoint
(\texttt{s3\_endpoint}) et de style d'adressage
(\texttt{s3\_url\_style='path'}), MinIO n'étant pas AWS S3. \\

\label{gloss:cosu}CoSu (Comité de suivi) &
Réunion mensuelle de pilotage entre le chef·fe de projet, le
commanditaire opérationnel (Karim BELAÏD) et la sponsor (Éléonore
RAKOTO). \\

\label{gloss:gantt}Diagramme de Gantt &
Représentation graphique d'un planning : une barre horizontale par
tâche, positionnée dans le temps, montrant enchaînements et
chevauchements entre phases -- voir
figure~\ref{fig:gantt}\footnote{Maintenu en Mermaid (texte versionné
dans le dépôt) plutôt que dans un outil graphique externe.}. \\

\label{gloss:bi}BI (\textit{Business Intelligence}) &
Ensemble des pratiques et outils visant l'\textbf{aide à la prise de
décision} à partir de données consolidées (tableaux de bord,
indicateurs, restitution visuelle) -- par opposition à la donnée brute
ou à son seul stockage. Chez Omega Logistics : l'entrepôt
\textbf{OMEGA BI} (bloc 3) et sa restitution dans \textbf{Grafana}
(C16bis). \\

\end{longtable}
